\documentclass[12pt]{article}

\usepackage[margin=0.75in]{geometry}
\usepackage{setspace}
\usepackage{parskip}
\usepackage{hyperref}
\usepackage{titlesec}
\usepackage{xcolor}
\usepackage{array}
\usepackage{ragged2e}

\definecolor{PrimaryRed}{HTML}{8B1E1E}
\definecolor{SoftBlue}{HTML}{DDEEFF}
\definecolor{DarkGray}{HTML}{333333}

\hypersetup{
    colorlinks=true,
    linkcolor=black,
    urlcolor=black,
    citecolor=black
}

\titleformat{\section}{\large\bfseries\color{PrimaryRed}}{}{0em}{}
\titleformat{\subsection}{\normalsize\bfseries}{\thesubsection}{0.5em}{}

\setstretch{1.15}

\begin{document}

\begin{titlepage}
\newgeometry{margin=0.85in}
\thispagestyle{empty}

\begin{center}
    \vspace*{0.25in}
    {\large University of Rochester\\}
    {\normalsize Warner School of Education\\[0.35in]}

    {\color{PrimaryRed}\rule{\textwidth}{1.4pt}}\\[0.22in]
    {\Huge\bfseries Joy-Oriented Literacy Pursuit\\[0.12in]}
    {\Large Formal Reflective Journal\\[0.22in]}
    {\color{PrimaryRed}\rule{0.72\textwidth}{0.7pt}}\\[0.32in]

    \colorbox{SoftBlue}{%
        \parbox{0.86\textwidth}{%
            \centering\vspace{0.12in}
            \textbf{Classroom Technology, Design, and Literacy as Joy}\\[0.05in]
            Kindergarten circuits, Grade 4 Minecraft, and Grade 5 airplane/fan-circuit design
            \vspace{0.12in}
        }%
    }
\end{center}

\vspace{0.75in}

\begin{center}
\renewcommand{\arraystretch}{1.35}
\begin{tabular}{>{\bfseries\RaggedLeft}p{1.55in} p{4.85in}}
Student & Piter Garcia \\
Course & EDU498: Culturally-Historically Responsive Literacies \\
Class Session & Class 4 \\
Assignment & Joy-Oriented Literacy Pursuit \\
Context & Teaching Placement Reflection \\
Submission Date & May 31, 2026 \\
\end{tabular}
\end{center}

\vfill

\begin{center}
\begin{minipage}{0.78\textwidth}
\centering
\itshape
Joy-oriented literacy is not extra. It is part of access. When students feel capable, connected, and allowed to participate as themselves, literacy becomes a tool for liberation, self-determination, and empowerment.
\end{minipage}
\end{center}

\restoregeometry
\end{titlepage}
\clearpage

\section*{What I Created and Who I Shared It With}

For my joy-oriented literacy pursuit, I created and facilitated classroom-based technology and design activities where students could experience literacy through building, collaboration, speaking, movement, problem-solving, and shared excitement. Instead of treating literacy only as reading or writing, I focused on moments where students used language, materials, technology, and relationships to make meaning. The main examples from my teaching placement were the Kindergarten circuit and fan lesson, the Grade 4 Minecraft activity, and the Grade 5 airplane and fan-circuit challenge. These activities were shared with students in my placement classrooms, but they were also shared with the classroom community itself, because students built, watched, helped, responded, and learned from one another in real time.

In the Kindergarten circuit lesson, students gathered around a demonstration where they learned about batteries, switches, wires, Play-Doh, and a fan by watching, naming parts, holding materials, connecting pieces, and counting down together before turning the fan on. The lesson became a shared literacy event because students were not only receiving information; they were participating physically, verbally, visually, and socially in making the concept meaningful. In Grade 4, students who finished their Minecraft work were invited to help classmates who were not done, which turned digital literacy into peer teaching and community responsibility. In Grade 5, students worked in teams to build airplane and fan circuits, with different roles such as building wings, shaping the body, making wheels, and connecting the fan. These moments helped me see that a literacy product does not always have to be a written artifact. Sometimes the product is a shared experience where students become makers, helpers, designers, and explainers.

\section*{My Creative and Learning Process}

My creative process began with a question: how can technology lessons become more than task completion? I wanted students to feel that literacy was something they could touch, build, speak, test, revise, and share. Instead of designing lessons where I explained everything and students simply copied the steps, I looked for ways to make students active contributors to the learning process. This meant creating roles, using materials, inviting peer support, and allowing students to experience success through collaboration rather than only through individual performance.

For the circuit and airplane activities, I used materials such as switches, batteries, wires, Play-Doh, fans, and model-building pieces to help students connect vocabulary to action. Students did not only hear the word ``switch''; they saw it, held it, connected it, and watched what happened when it worked. In the Grade 5 activity, students had to decide who would make each part of the airplane, which made the activity collaborative and distributed. In the Grade 4 Minecraft lesson, I used the students who finished early as resources for the class instead of treating completion as the end of learning. That shifted the process from ``I am done'' to ``I can help someone else succeed.''

\newpage
This process also taught me to notice joy as part of learning design. Joy did not happen only because the activities were fun. It happened because students had access to materials, roles, language, peer support, and visible success. The creative process helped me understand that joy has to be planned for with intention. Students need enough structure to feel safe and enough freedom to feel ownership. When those conditions exist together, learning becomes more than compliance; it becomes participation.

\section*{How This Cultivated Joy, Love, and Aesthetic Fulfillment}

This literacy pursuit cultivated joy because students were able to see their learning become real. In the Kindergarten circuit lesson, the joy came from collective anticipation. Students watched the pieces connect, counted down together, and saw the fan turn on. That moment mattered because the science was not abstract anymore. It became visible, audible, physical, and shared. The joy was not only in the fan working; it was in the class experiencing the success together.

The pursuit also cultivated love through relationships. In Grade 4 Minecraft, students who already understood the task helped classmates who were still working. That kind of peer support changed the emotional meaning of the activity. Success was not private or competitive; it became something students could give to each other. In Grade 5, the airplane and fan-circuit challenge created aesthetic fulfillment because students were building something with shape, motion, design, and purpose. Their work was not only ``correct'' or ``incorrect.'' It had form, creativity, visible effort, and a relationship to the physical world.

These activities reminded me that joy in literacy does not always look quiet or neat. Sometimes joy looks like a student trying again, a group arguing productively over what goes where, a fan finally turning on, a peer stepping in to help, or a student realizing that their hands, ideas, and voice matter. For me, that is aesthetic fulfillment: the moment when learning feels alive, when students can see themselves inside the work, and when the classroom becomes a place where knowledge is built rather than simply delivered.

\section*{Muhammad's Literacy Pursuit}

This project connects to Muhammad's idea of literacy pursuit because the activities were not designed only to teach isolated skills. They created opportunities for students to build identity, skills, intellect, criticality, and joy. Students were developing vocabulary and technical understanding, but they were also learning how to see themselves as builders, helpers, problem-solvers, and contributors. That matters because literacy should not only prepare students to answer questions. Literacy should help students understand themselves, participate in community, and experience agency.

\newpage
Muhammad's framework helps me understand that joy has to be designed for, not assumed. If I had taught the circuit lesson only as a teacher demonstration, students might have understood the concept, but they would not have owned it in the same way. By allowing students to touch materials, contribute to the build, name what they saw, and participate in the countdown, the lesson became a literacy pursuit instead of a simple activity. The same is true for Minecraft and the airplane challenge. Students were reading instructions, interpreting systems, using digital and physical tools, communicating with classmates, and making meaning through action.

This also connects to culturally and historically responsive literacy because it expands what counts as literacy. Students bring different bodies, languages, identities, cultures, learning styles, and access needs into the classroom. A joy-oriented literacy pursuit gives students more than one way to enter the work. It recognizes that literacy can include collaboration, design, movement, translation, troubleshooting, visual thinking, and oral explanation. In that way, joy is not separate from rigor. Joy becomes part of how students access rigor.

\section*{How This Informs My Future Teaching}

This experience will inform my teaching by reminding me that literacy is not limited to written language. Literacy can be building, coding, designing, explaining, translating, collaborating, troubleshooting, helping, and trying again. In my future classroom, I want to design activities where students have multiple entry points and multiple ways to show understanding. This is especially important for students who may be multilingual, neurodivergent, disabled, chronically ill, culturally diverse, or simply different from the ``default'' learner schools often imagine.

In practice, I would continue designing projects where students create something meaningful and share it with an authentic audience. Students could build Minecraft models connected to a science concept, create bilingual or multimodal explanations of a design challenge, make collaborative classroom murals, create technology tutorials for peers, or build simple circuits and explain how they work through speaking, drawing, writing, or video. I would also make peer teaching part of the classroom culture, because students often understand each other's barriers in ways adults miss.

My biggest takeaway is that joy-oriented literacy is not extra. It is part of access. When students feel capable, connected, and allowed to participate as themselves, literacy becomes a tool for liberation, self-determination, and empowerment. This pursuit helped me see that my teaching practice should not only ask whether students completed the task. It should also ask whether students had the chance to feel ownership, beauty, relationship, and possibility inside the learning.

\section*{Evidence Notes}

This reflection draws on placement evidence from the Kindergarten circuit and fan lesson, the Grade 4 Minecraft peer-support activity, and the Grade 5 airplane and fan-circuit challenge. These examples show students using literacy through oral language, digital tools, physical materials, collaboration, design, and shared problem-solving.

\end{document}
